
\documentclass{report}

\input{preamble}
\input{macros}
\input{letterfonts}

\title{\Huge{Programmazione Web}\\Sistemi Informativi Aziendali}
\author{\huge{Caberlotto Giovanni}}
\date{}

\begin{document}

\maketitle
\newpage% or \cleardoublepage
% \pdfbookmark[<level>]{<title>}{<dest>}
\pdfbookmark[section]{\contentsname}{toc}
\tableofcontents
\pagebreak

\chapter{Come funziona il internet} % (fold)
\label{chap:Come funziona il internet}

% chapter Come funziona il internet (end)

\chapter{Programmazione lato client} % (fold)
\label{chap:Programmazione lato client}

\section{HTML}

\subsection{Che cos'è l'HTML}
\\
HTML è l'acronimo di HyperText Markup Language ("Linguaggio di contrassegno per gli Ipertesti") e non è un linguaggio di programmazione. Si tratta invece di un linguaggio di markup (di 'contrassegno' o 'di marcatura'), che permette di indicare come disporre gli elementi all'interno di una pagina.
\\
\begin{lstlisting}[language=HTML, caption=Esempio di codice in HTML]
<!DOCTYPE html>
<html lang="en">
  <head>
    <title></title>
    <meta charset="UTF-8">
    <meta name="viewport" content="width=device-width, initial-scale=1">
    <link href="css/style.css" rel="stylesheet">
  </head>
  <body>
  
  </body>
</html>
\end{lstlisting}
% chapter Programmazione lato client (end)
\\
Queste indicazioni vengono date attraverso degli appositi marcatori, detti tag ('etichette'), che hanno la caratteristica di essere inclusi tra parentesi angolari (ad es, <img> è il segnaposto di un'immagine).
\\
\\
Con HTML quindi indichiamo, attraverso i tag, quali elementi dovranno apparire su uno schermo e come essi debbano essere disposti. Tutte queste indicazioni sono contenute in un documento HTML, spesso detto "Pagina HTML". Una pagina HTML è rappresentata da un file di testo, ovvero un file che possiamo modificare con programmi come notepad e in genere hanno un nome che finisce con l'estensione .html.
\\
\\
Ciò significa che HTML non possiede i costrutti propri della programmazione, come i meccanismi "condizionali", che consentono di reagire in modo diverso a seconda del verificarsi di una condizione ("in questa situazione fai questo, in quest'altra fai quest'altro"), o i costrutti iterativi ("ripeti questa azione, finché non succede questo").
\\
\nt{Si parla in questo caso di linguaggio dichiarativo, che serve appunto ad indicare cosa deve apparire sullo schermo (testi, immagini, suoni, etc.), dove e in che sequenza. Nel caso di linguaggi in cui specifichiamo algoritmi precisi con "strutture di controllo", come C, C++, Java, o anche PHP e JavaScript, parliamo di "linguaggi imperativi".}
\\
\\

\end{document}
